% sintaxismark-doc-es.tex
% Documentación en español para sintaxismark.sty v0.6.21
% Compilar con LuaLaTeX. Ejecutar dos veces porque TikZ usa remember picture.

\documentclass[11pt,a4paper]{article}

\usepackage{fontspec}
\usepackage[spanish]{babel}
\usepackage{geometry}
\geometry{margin=2.5cm}
\usepackage{booktabs}
\usepackage{longtable}
\usepackage{xcolor}
\usepackage{listings}
\usepackage{hyperref}
\usepackage{parskip}
\usepackage{microtype}
\usepackage{gb4e}\noautomath

\usepackage{sintaxismark}

\lstset{
  language=[LaTeX]TeX,
  basicstyle=\ttfamily\small,
  columns=fullflexible,
  breaklines=true,
  frame=single,
  backgroundcolor=\color{gray!8},
  keywordstyle=\color{blue!70!black},
  commentstyle=\color{gray!70!black}
}

\title{\texttt{sintaxismark}: documentación}
\author{}
\date{Versión 0.6.21 --- 2026-06-15}

\begin{document}
\maketitle
\tableofcontents

\section{Introducción}

\texttt{sintaxismark} es un paquete de LaTeX para marcar constituyentes sintácticos, funciones gramaticales, núcleos y tramos directamente sobre texto corrido. Está pensado para materiales de gramática, ejemplos de clase, diapositivas y análisis lingüísticos compactos en los que un árbol completo sería innecesario o visualmente demasiado pesado.

El paquete usa superposiciones de TikZ. Permite dibujar corchetes inferiores abiertos, líneas superiores de alcance, etiquetas flotantes, subrayados, flechas, marcas anidadas y cortes manuales de tramos largos.

\section{Dependencias}

El paquete carga o requiere los siguientes componentes:

\begin{itemize}
  \item \texttt{tikz}, con las bibliotecas \texttt{calc} y \texttt{arrows.meta};
  \item \texttt{xparse};
  \item \texttt{xcolor};
  \item \texttt{expl3};
  \item \texttt{varwidth}.
\end{itemize}

El paquete funciona bien con \texttt{gb4e}, aunque \texttt{gb4e} no es obligatorio salvo que se quieran ejemplos lingüísticos numerados.

Como el paquete usa el mecanismo \texttt{remember picture} de TikZ, normalmente conviene compilar dos veces.

\section{Carga del paquete}

Coloque \texttt{sintaxismark.sty} en el mismo directorio que el documento o instálelo en el árbol local de TeX. Luego cárguelo en el preámbulo:

\begin{lstlisting}
\usepackage{sintaxismark}
\end{lstlisting}

También se puede seleccionar un esquema visual al cargar el paquete:

\begin{lstlisting}
\usepackage[docencia]{sintaxismark}
\usepackage[impresion]{sintaxismark}
\usepackage[diapositivas]{sintaxismark}
\usepackage[publicacion]{sintaxismark}
\end{lstlisting}

El esquema activo también puede modificarse dentro del documento:

\begin{lstlisting}
\setsintaxisschema{diapositivas}
\end{lstlisting}

\section{Ejemplo mínimo}

\begin{lstlisting}
\documentclass{article}
\usepackage{gb4e}\noautomath
\usepackage[docencia]{sintaxismark}

\begin{document}

\begin{exe}
\ex \SES{Juan} \PVS{\NV{compró} \OD{manzanas}}.
\end{exe}

\end{document}
\end{lstlisting}

\section{Comandos básicos}

\subsection{Corchetes inferiores abiertos: \texttt{\textbackslash openbox}}

\begin{lstlisting}
\openbox{OD}{Juan compró manzanas}
\end{lstlisting}

El primer argumento es la etiqueta. El segundo argumento es el texto anotado. Se pueden pasar claves opcionales entre corchetes:

\begin{lstlisting}
\openbox[color=blue,linew=1pt,boxd=4ex]{OD}{manzanas}
\end{lstlisting}

\subsection{Líneas superiores de alcance: \texttt{\textbackslash spanline}}

\begin{lstlisting}
\spanline{SES}{Juan} \spanline{PVS}{compró manzanas}
\end{lstlisting}

Este comando dibuja una línea horizontal superior con extremos verticales y una etiqueta sobre el tramo.

\subsection{Etiquetas flotantes: \texttt{\textbackslash tagabove} y \texttt{\textbackslash tagbelow}}

\begin{lstlisting}
\tagabove{ST}{Juan compró manzanas}
\tagbelow{NN}{Juan}
\end{lstlisting}

\subsection{Subrayado: \texttt{\textbackslash synul}}

\begin{lstlisting}
\synul{compró}
\end{lstlisting}

Sirve para marcar núcleos u otros elementos destacados.

\section{Atajos predefinidos}

El paquete define atajos para etiquetas frecuentes de análisis sintáctico y de gramática escolar.

\begin{longtable}{ll}
\toprule
Comando & Etiqueta por defecto \\
\midrule
\verb|\SES{...}| & \textsc{ses} \\
\verb|\ST{...}| & \textsc{st} \\
\verb|\SEC{...}| & \textsc{sec} \\
\verb|\PVS{...}| & \textsc{pvs} \\
\verb|\PVC{...}| & \textsc{pvc} \\
\verb|\OS{...}| & \textsc{os} \\
\verb|\OD{...}| & \textsc{od}, con subrayado \\
\verb|\OI{...}| & \textsc{oi}, con subrayado \\
\verb|\MD{...}| & \textsc{md} \\
\verb|\MI{...}| & \textsc{mi} \\
\verb|\NN{...}| & \textsc{nn}, con subrayado \\
\verb|\NV{...}| & \textsc{nv}, con subrayado \\
\verb|\NP{...}| & \textsc{np} \\
\verb|\NA{...}| & \textsc{na} \\
\verb|\NAD{...}| & \textsc{nadv} \\
\bottomrule
\end{longtable}

Otros atajos duales son \verb|\CCL|, \verb|\COMP|, \verb|\CAg|, \verb|\CREG|, \verb|\CSim|, \verb|\CL|, \verb|\CM|, \verb|\PSO|, \verb|\PSNO|, \verb|\POO|, \verb|\PONO|, \verb|\ACT|, \verb|\ACL|, \verb|\ACM|, \verb|\ACF|, \verb|\ACC|, \verb|\ACCant|, \verb|\ACI|, \verb|\ACFin|, \verb|\MG|, \verb|\Ap|, \verb|\MF|, \verb|\IPCR| e \verb|\Incl|.

\section{Formas con y sin estrella}

Desde la versión 0.6.19, muchos atajos tienen un comportamiento dual:

\begin{itemize}
  \item la forma sin estrella usa una etiqueta inferior flotante, normalmente mediante \verb|\tagbelow|;
  \item la forma con estrella usa un corchete inferior abierto, mediante \verb|\openbox|.
\end{itemize}

\begin{lstlisting}
\OD{manzanas}     % etiqueta inferior + subrayado
\OD*{manzanas}    % corchete inferior abierto

\MI{de lino}      % etiqueta inferior
\MI*{de lino}     % corchete inferior abierto
\end{lstlisting}

\section{Esquemas visuales}

El paquete incluye cuatro esquemas:

\begin{longtable}{ll}
\toprule
Esquema & Uso previsto \\
\midrule
\texttt{docencia} & materiales de clase y ejercicios \\
\texttt{impresion} & apuntes impresos compactos \\
\texttt{diapositivas} & diapositivas y material proyectado \\
\texttt{publicacion} & salida más compacta orientada a publicación \\
\bottomrule
\end{longtable}

\section{Claves gráficas}

La mayoría de los comandos aceptan opciones en formato clave-valor:

\begin{lstlisting}
\openbox[color=blue,linew=1pt,boxd=4ex]{OD}{manzanas}
\spanline[color=red,liney=3ex]{SES}{Juan}
\tagbelow[tagx=0.5ex,tagy=2ex]{MD}{el}
\end{lstlisting}

\begin{longtable}{ll}
\toprule
Clave & Función \\
\midrule
\texttt{color} & color de la anotación \\
\texttt{linew} & grosor de línea \\
\texttt{liney} & posición vertical de las líneas superiores \\
\texttt{linegap} & separación entre texto y extremos superiores \\
\texttt{labsep} & separación entre línea superior y etiqueta \\
\texttt{boxd} & profundidad del corchete inferior \\
\texttt{boxsep} & separación entre corchete inferior y etiqueta \\
\texttt{labelsep} & separación general de etiquetas \\
\texttt{labelsepboxes} & separación de etiquetas para cajas \\
\texttt{nestedgap} & separación extra para marcas anidadas \\
\texttt{labheight} & altura estimada de etiqueta para cálculos de espaciado \\
\texttt{tagy} & posición vertical de etiquetas flotantes \\
\texttt{tagx} & desplazamiento horizontal de etiquetas flotantes \\
\texttt{underw} & grosor del subrayado \\
\texttt{labelalign} & alineación de etiqueta: \texttt{left}, \texttt{center} o \texttt{right} \\
\texttt{leftend}, \texttt{rightend} & mostrar u ocultar extremos verticales \\
\texttt{leftarrow}, \texttt{rightarrow} & agregar flechas \\
\texttt{arrowlen} & longitud de punta de flecha \\
\texttt{arrowtip} & tipo de punta TikZ \\
\texttt{arrowscale} & escala de punta de flecha \\
\texttt{arrowinset} & inserción de punta de flecha \\
\bottomrule
\end{longtable}

\section{Ejemplos largos y saltos de línea}

Para ejemplos que deben partirse en más de una línea, use \verb|\SMlinewrap|:

\begin{lstlisting}
\begin{exe}
\ex \SMlinewrap{El \MI{departamento de unos monoblocks deteriorados de tres pisos de color bordó} quedaba en la planta baja.}
\end{exe}
\end{lstlisting}

Si un tramo no debe partirse entre líneas, use las variantes sin corte:

\begin{lstlisting}
\openboxNB{MI}{departamento de unos monoblocks deteriorados}
\spanlineNB{SES}{El departamento de unos monoblocks deteriorados}
\end{lstlisting}

\section{Tramos partidos manualmente}

El comando \verb|\breakspan| sirve cuando un tramo debe dividirse manualmente en dos partes visuales. El primer argumento se muestra como tramo cortado hacia la derecha; el segundo, como tramo cortado hacia la izquierda.

\begin{lstlisting}
\breakspan
  {\SES{\MI{El departamento de unos monoblocks}}}
  {\SES{\MI{quedaba en la planta baja}}}
\end{lstlisting}

\section{Notas prácticas}

\begin{itemize}
  \item Compile dos veces si las etiquetas o líneas aparecen mal ubicadas en la primera pasada.
  \item Al usar \texttt{gb4e}, cargue \texttt{gb4e} antes que \texttt{sintaxismark}.
  \item Para ejemplos largos, use \verb|\SMlinewrap|, \verb|\openboxNB|, \verb|\spanlineNB| o \verb|\breakspan|.
  \item Las etiquetas anidadas se espacian automáticamente mediante propagación de profundidad, mejorada en la versión 0.6.21.
\end{itemize}

\section{Notas de versión}

La versión 0.6.21 mejora el espaciado de \verb|\tagbelow| anidados mediante encadenamiento de profundidad, de modo análogo a \verb|\openbox|. Esto evita que las etiquetas inferiores anidadas se superpongan entre sí o con etiquetas de la línea siguiente.

\end{document}
